\documentclass[dvipsnames,draft]{beamerikz2}
% MODES: ready, final, draft, short, brief
% FONTS: small, basic, large
% other: plain

% packages --------------------------------------------------------------------

% \usepackage{}
\usepackage{bbold}
\usepackage{varwidth}
\usepackage{mathdots}
\usepackage{pgfplots}
\usepackage{svg}
\usetikzlibrary{arrows,automata}
\usetikzlibrary{decorations.pathmorphing,shapes}
% document data ---------------------------------------------------------------

\sBzTitle{Nonlinear extensions of linear recurrence sequences}

\sBzAuthor{Filip Mazowiecki}

% macros ----------------------------------------------------------------------

% \newcommand{\w}{\omega}

% =============================================================================

%eg ie
\newcommand*{\eg}{e.g.\@\xspace}
\newcommand*{\ie}{i.e.\@\xspace}

%Automata
\newcommand{\Aa}{\mathcal{A}}
\newcommand{\Bb}{\mathcal{B}}
\newcommand{\Cc}{\mathcal{C}}
\newcommand{\X}{\mathcal{X}}
\newcommand{\Mm}{\mathcal{M}}
\newcommand{\Oo}{\mathcal{O}}

%VECTORS
\newcommand{\uu}{\boldsymbol{u}}
\newcommand{\vv}{\boldsymbol{v}}
\newcommand{\aaa}{\boldsymbol{a}}
\newcommand{\bbb}{\boldsymbol{b}}

%SETS
\newcommand{\Z}{\mathbb{Z}}
\newcommand{\N}{\mathbb{N}}
\newcommand{\Q}{\mathbb{Q}}
\newcommand{\R}{\mathbb{R}}
\newcommand{\Qpos}{\mathbb{Q}_{\ge 0}}
\newcommand{\bbS}{\mathbb{S}}
\renewcommand{\S}{\bbS}

%stuff
\newcommand{\wh}[1]{\widehat{#1}}
\newcommand{\set}[1]{\{ #1 \}}
\newcommand{\sem}[1]{\left\llbracket #1 \right\rrbracket}
\newcommand{\transpose}{\mathsf{T}}
\newcommand{\td}[1]{\tilde{#1}}

%semirings
\newcommand{\semiringS}{\S(\oplus,\odot)}
\newcommand{\semiringQ}{\Q(+,\cdot)}
\newcommand{\semiringQm}{\Q(\max,\cdot)}
\newcommand{\semiringZ}{\Z(\min,+)}
\newcommand{\semiringZR}{\R(\min,+)}
\newcommand{\semiringN}{\N(\min,+)}
\newcommand{\zero}{\mathbb{0}}
\newcommand{\one}{\mathbb{1}}

\newcommand{\poly}{\operatorname{poly}}

%fill
\newcommand{\background}[2]{
\dBzH[0.5]
\path[fill=#1,rounded corners=3pt] (-10,\bzH) rectangle (10,\bzH-#2);
\iBzH[0.5]
}

\begin{document}

\begin{bzPlainFrame}%[show]
\bzOn{
	\iBzH[2.5]
	\bzCenter[bzEB, scale=1.4]{Nonlinear extensions of}
	\iBzH[0.6]
	\bzCenter[bzEB, scale=1.4]{linear recurrence sequences}

	\iBzH[2.2]
	\bzCenter{Filip Mazowiecki}
	\iBzH[0.4]
	\bzCenter[scale=0.7]{\sc Uniwersytet Warszawski}

	\iBzH[1.8]
	\bzCenter{SAMSA 2026}

}
\end{bzPlainFrame}

%LRS
\begin{bzFrame}%[show]
\bzOn{
	\bzCenter[bzEB]{Linear recurrence sequences (LRS)}
\bzText{$a : \N \to \Q$}
	\iBzH[0.5]
	\bzText{For example {\textcolor{bzRed}{$a_n = n^2$}}}

	\iBzH[2]
	
    \bzCenter[align=left]{$\textcolor{bzRed}{a_0} = 0$,\quad $\textcolor{bzRed}{a_{n+1}} = \textcolor{bzRed}{a_n} + 2\textcolor{bzBlue}{b_n} + \textcolor{bzGreen}{c_n}$ \\
    $\textcolor{bzBlue}{b_0} = 0$,\quad $\textcolor{bzBlue}{b_{n+1}} = \textcolor{bzBlue}{b_n} + \textcolor{bzGreen}{c_n}$ \\
    $\textcolor{bzGreen}{c_0} = 1$,\quad $\textcolor{bzGreen}{c_{n+1}} = \textcolor{bzGreen}{c_n}$}
}

\iBzH[2]

\bzOn{
\bzCenter{$\textcolor{bzRed}{a_n} =
\begin{pmatrix}
0 & 0 &1
\end{pmatrix}
\begin{pmatrix}
1 & 0 & 0\\
2 & 1 & 0\\
1 & 1 & 1
\end{pmatrix}^n
\begin{pmatrix}
1\\
0\\
0
\end{pmatrix}
$}

\iBzH

\bzText{In general $u_n = I^{\transpose}M^nF$}
}

\iBzH[1]

\bzOn{

    \bzText{Can be also defined by  {\textcolor{bzRed}{$a_{n+3} = 3a_{n+2} - 3a_{n+1} + a_n$}}}
    \bzText{In general $u_{n+k} = \sum_{i=0}^{k-1} \alpha_i u_{n+i}$}
}
\end{bzFrame}

%LRS PROPERTIES
\begin{bzFrame}
\bzOn{
\bzCenter[bzEB]{Many nice properties of LRS}

\iBzH[1]

\bzItem{Unique, computable minimal degree}
\bzText[bzRed]{($d = 3$ for $a_n = n^2$)}
}

\iBzH[1]

\bzOn{
\bzItem{Closure properties, given LRS $a_n$, $b_n$:}
\bzText[bzRed]{$c_n = a_n + b_n$, \quad $c_n = a_n \cdot b_n$, \quad $c_n = a_n * b_n = \sum_{i<n}a_i \cdot b_{n-i-1}$}
}

\iBzH[1]

\bzOn{
\bzItem{Closed form}
\bzText[bzRed]{$a_n = \sum_i p_i(n) \cdot \lambda_i^n$}
\bzText{(for $n$ big enough)}
}

\iBzH[1]

\bzOn{
\bzText{e.g. $F_n = \frac{1}{\sqrt{5}}\varphi^n + \frac{1}{\sqrt{5}} \psi^n$}
}
\end{bzFrame}

%WEIGHTED
\begin{bzFrame}
\bzOn{
\bzCenter[bzEB]{Weighted automata over the rational field}

\iBzH[1]

\bzItem{$\Aa$ is defined by $\Sigma = \set{a,b,c}$, \quad $M_a,M_b,M_c \in \Q^{d \times d}, \quad I,F \in \Q^d$}
}

\iBzH[1]

\bzOn{
\bzText[bzRed]{$\Aa : \Sigma^* \to \Q$}
\bzText[bzRed]{$\Aa(w) = I^{\transpose} M_w F$, where $M_w = M_{w_1}\ldots M_{w_n}$ for $w = w_1\ldots w_n$.}
}

\iBzH[1]

\bzOn{
\bzItem{LRS are weighted automata for $|\Sigma| = 1$.}
}

\iBzH[1]

\bzOn{
\bzItem{We consider extensions of LRS from this point of view.}
}

\bzOn{
\bzText{We work over $\Q$}
}
\end{bzFrame}

%Polynomial
\begin{bzFrame}
\bzOn{
\bzCenter[bzEB]{Polynomial sequences (polyrec)}

\iBzH[1]

\bzText{Change matrices $M_a$ to polynomial maps}
}

\bzOn{
\bzText{Polynomial automata, cost-register automata, polynomial recurrent relations}
\bzText{\textcolor{bzGray}{[Benedikt et al., 2017], [Alur et al., 2013], [S\'{e}nizergues, 2007]}}
}


\iBzH[1]

\bzOn{
\bzText{Example $a_n = 2^{2^n}$}
\bzText[bzRed]{$a_0 = 2$, \quad $a_{n+1} = a_n^2$, \quad $f_1(x) =  x^2$}
}

\iBzH[1]

\bzOn{
\bzText{Example $a_n = n!$}
\bzText[bzRed]{$a_0 = 1$, \quad $a_{n+1} = a_n \cdot b_n$, \quad $f_1(x,y) = xy$}
\bzText[bzRed]{$b_0 = 1$, \quad $b_{n+1} = b_n + 1$, \quad $f_2(x,y) = y+1$}
}

\iBzH[1]

\bzOn{
\bzText{Example $a_n = 2^{2^n} + n!$, \quad or $a_n = 2^{2^n} \cdot n!$}
}
\end{bzFrame}

%holonomic
\begin{bzFrame}
\bzOn{
\bzCenter[bzEB]{Holonomic/P-recursive sequences}

\iBzH[1]

\bzText{In the automata definition generalise to context-free grammars (CFGs)}
\bzText{Weighted/probabilistic CFGs \textcolor{bzGray}{[Chomsky and Sch\"{u}tzenberger, 1959]}}
}

\iBzH[1]

\bzOn{
\bzText{Equivalently: allow polynomials (in $n$) in the recursion}
}

\iBzH[1]

\bzOn{
\bzText{Example $a_n = n!$}
\bzText[bzRed]{$a_0 = 1$, \quad $a_{n+1} = a_n \cdot n$}
}

\iBzH[1]

\bzOn{
\bzText{Example $C_n$ Catalan numbers}
\bzText[bzRed]{$C_0 = 1$, \quad $(n+2)C_{n+1} = (4n+2)C_n$}
}

\iBzH[1]

\bzOn{
\bzText{Example $a_n = C_n + n!$, \quad $a_n = C_n \cdot n!$, \quad $a_n = C_n * n!$}
}
\end{bzFrame}

%wmso
\begin{bzFrame}
\bzOn{
\bzCenter[bzEB]{MSO sequences}

\iBzH[1]

\bzText{Present in the logic framework, replace quantifiers with aggregators}
\bzText{Weighted MSO \textcolor{bzGray}{[Droste and Gastin, 2005]}}
}

\iBzH[1]

\bzOn{
\bzText{MSO formula evaluate to $0$ or $1$, \quad constants evaluate to themselves,}
\bzText{$\wedge$, $\vee$ become $\cdot$ and $+$, \quad $\exists$,  $\forall$ become $\sum$ and $\prod$}
}

\iBzH[1]

\bzOn{
\bzText{Example $a_n = n^n$ and $b_n = n!$}
\bzText[bzRed]{$a_n = \prod_{x} \sum_y 1$, \quad $b_n = \prod_x \sum_y [x \le y]$}
}

\iBzH[1]

\bzOn{
\bzText{Example $a_n = 2^n$ and $b_n = 2^{2^n}$}
\bzText[bzRed]{$a_n = \sum_X 1$, \quad $b_n = \prod_X 1$}
}
\end{bzFrame}

%expressivness 
\begin{bzFrame}
\bzOn{
\bzCenter[bzEB]{Three nonlinear extensions}


}
\end{bzFrame}


%asymptotics
\begin{bzFrame}
\bzOn{
\bzCenter[bzEB]{Asymptotics of sequences}

\iBzH[1]

\bzText{Bound on $|a_n| \le f(n)$?}
}
\end{bzFrame}


%CONCLUSIONS
\begin{bzFrame}%[show]
\bzOn{
\bzCenter[bzEB]{Conclusion}

	\iBzH[1]

\bzItem{Boundedness proof is inspired by {\color{gray}[Simon 1994]}}
	\iBzH[1]
\bzText{we believe it should work in more dimensions}
}

\end{bzFrame}

% =============================================================================

\end{document}
